% Author: Simon-Pierre Boucher — contact@spboucher.ai
%
% ═══════════════════════════════════════════════════════════════════════
% 7. DISCUSSION
% ═══════════════════════════════════════════════════════════════════════
\section{Discussion}
\label{sec:discussion}

The results support four main messages, each associated with a specific modeling framework.

\subsection{Message 1: OLS Remains Useful for Interpretable Associations}

The semi-log OLS model explains 63.4\% of the variation in log listing prices using 62 regressors, comparable to the typical hedonic $R^2$ in the literature \citep{sirmans2005composition}. Its primary value lies in interpretability: the coefficients provide conditional mean associations in a transparent functional form. Regularized linear models (Ridge, Lasso, Elastic Net) achieve essentially identical performance ($R^2 = 0.628$--$0.630$), confirming that the OLS specification is not overfit and that the 20-percentage-point gap relative to XGBoost reflects genuine non-linearities and interactions, not parametric overfitting.

The progression from region FE ($R^2 = 0.634$) to state FE ($0.678$) to ZIP3 FE ($0.725$) reveals that a large portion of the unexplained variation is spatial in nature. The ZIP3 FE model closes roughly half the gap between baseline OLS and XGBoost, suggesting that fine-grained geographic controls---rather than non-linear functional forms---account for much of the ML advantage. However, ZIP3 FE introduces 886 parameters and sacrifices the parsimony that makes OLS attractive for economic interpretation.

The imputation sensitivity analysis reveals that the baseline lot-size elasticity (0.02) is substantially attenuated by state-median imputation; the complete-case estimate (0.06) is more plausible but applies to a selected subsample. This finding underscores the importance of data quality audits in hedonic research.

\subsection{Message 2: Quantile Regression Adds Distributional Insight}

Quantile regression reveals that OLS coefficients mask substantial distributional heterogeneity. The garage gradient ratio (10.3:1 between $\tau = 0.10$ and $\tau = 0.90$) and the pool gradient reversal (insignificant at lower quantiles, strongly positive at upper quantiles) demonstrate that the same attribute can play fundamentally different roles across market segments. Inter-quantile Wald tests confirm that these differences are statistically significant for 11 of 13 tested variables, with garage ($z = 28.92$) being the most significantly heterogeneous.

This has implications for housing policy: interventions that improve basic amenities (garages, heating systems) may generate the largest listing-price differentials at the lower end, while luxury amenity provision primarily differentiates upper-market properties.

The subsample stability analysis shows that most quantile coefficients are highly robust (CV $< 8\%$), but age stands out as unstable (CV = 78.6\%, sign stability = 80\%), reflecting genuine geographic heterogeneity in how property age relates to listing prices.

\subsection{Message 3: ML Performance Depends Critically on Validation Design}

The XGBoost model achieves substantially better predictions than OLS under random validation ($R^2 = 0.833$ vs.\ $0.630$), capturing non-linearities and interactions that the parametric specification cannot. However, the geographic holdout results provide an important corrective: under state-level holdout, XGBoost achieves only $R^2 = 0.425$--$0.547$ depending on specification.

The ablation analysis pinpoints what drives the ML gain. Neighborhood features contribute $+17.6$ pp in random $R^2$---the largest single gain---but only $+3.5$ pp in geographic holdout $R^2$. This disparity suggests that neighborhood variables (school ratings, Walk Score, etc.) are highly informative within observed markets but carry location-specific scale and meaning that may not transfer. The full model with region dummies actually \emph{reduces} geographic holdout $R^2$ by 7.6 pp relative to the market-status-only specification.

The most striking finding is that removing all geographic features \emph{improves} geographic holdout $R^2$ from 0.425 to 0.519. This is the opposite of what intuition might suggest and has direct implications for AVM deployment: in contexts requiring geographic generalization, simpler models without explicit geography may outperform richer ones.

\subsection{Message 4: SHAP Interprets Prediction, Not Economics}

SHAP values provide useful prediction-level decompositions that identify which features drive the XGBoost model's predictions for individual properties. The cross-model stability analysis ($\rho = 0.89$--$0.99$ across three tree models) addresses the common criticism that SHAP rankings are model-specific: while individual SHAP values differ, the ranking of feature importance is remarkably consistent, with the same six features dominating all three models.

However, SHAP values should not be equated with hedonic listing-price gradients:
\begin{itemize}[nosep]
    \item SHAP values decompose predictions, not the data-generating process. A feature with a large SHAP value may be predictively important because it proxies for unobserved variables, not because it has a large association with listing prices in the structural sense.
    \item Feature correlation distributes SHAP contributions among correlated features in ways that may not reflect economic importance.
    \item SHAP importance rankings differ from OLS coefficient rankings (e.g., school rating ranks 3rd by SHAP but has a counterintuitive negative OLS sign), reflecting the different questions each framework answers.
\end{itemize}

In the terminology of \citet{mullainathan2017machine}, SHAP is a tool for interpreting predictions. The hedonic gradient $\partial P / \partial z_k$ is a tool for understanding market structure. These serve different purposes and should not be conflated.
